\chapter{Evidence Ladder}
\label{app:evidence_ladder}

VLA 논문은 peer-reviewed paper, arXiv preprint, project page, company report, closed product announcement가 섞여 있다. 같은 문단에서 모두 인용할 수는 있지만, 같은 신뢰도로 읽으면 안 된다. 이 책은 다음 evidence ladder를 사용한다.

\begin{table}[h]
  \centering
  \caption{Evidence level used in this book.}
  \label{tab:evidence_ladder}
  \begin{tabularx}{\textwidth}{>{\bfseries}p{16mm}p{44mm}X}
    \toprule
    Level & Meaning & How to use it \\
    \midrule
    A & Peer-reviewed 또는 accepted paper이고, code/benchmark/real-robot evidence 중 적어도 하나가 비교적 재현 가능하다. & 본문 claim의 강한 anchor로 쓸 수 있지만, controller, split, hardware 조건은 여전히 확인한다. \\
    B & arXiv preprint이거나 project artifact가 있으며, 실험 protocol과 일부 artifact가 공개되어 있다. & Emerging result의 주요 근거로 쓸 수 있으나, peer review 전 claim과 artifact completeness를 구분한다. \\
    C & Project page, company report, closed model report, demo video 중심 source이다. & Industrial trend 또는 product direction을 설명하는 데 쓰되, 검증된 학술 evidence처럼 쓰지 않는다. \\
    D & Survey-only, blog-only, claim-level source이거나 원문 artifact 접근이 제한적이다. & 배경 신호로만 사용하고, 핵심 성능 claim의 직접 근거로 쓰지 않는다. \\
    \bottomrule
  \end{tabularx}
\end{table}

\section*{Company claim handling}

Closed industrial VLA는 field direction을 이해하는 데 중요하다. 그러나 Helix, Gemini Robotics, Gemini Robotics-ER 같은 source는 public reproducibility가 제한될 수 있다. 따라서 본문에서는 다음 규칙을 따른다.

\begin{enumerate}
  \item Company report는 ``관측되는 산업 방향''의 근거로 쓰고, 일반 SOTA 성능의 확정 근거로 쓰지 않는다.
  \item 공개되지 않은 model weight, data, controller, safety layer는 ``확인되지 않음''으로 남긴다.
  \item Demo success는 benchmark success와 분리한다.
  \item Closed VLA action model과 embodied reasoning model은 versioned naming으로 구분한다.
\end{enumerate}

이 ladder는 reviewer가 가장 먼저 묻는 질문, 즉 ``이 주장의 근거가 무엇인가''에 답하기 위한 장치이다.
